% Archivo generado por herramientas/generar_cifras.py a partir de herramientas/cifras_mapa.csv.
% No editar a mano: cada cifra se lee de data/results/ (regla 3 de CLAUDE.md).
% Para regenerarlo: python3 herramientas/generar_cifras.py

% Meses de la vía en niveles | data/results/05_dimensiones.csv | via=niveles | n
\newcommand{\nMesesNiveles}{228}
% Meses de la vía transformada | data/results/05_dimensiones.csv | via=transformadas | n
\newcommand{\nMesesTransformadas}{227}
% Series del análisis (D1) | data/results/05_dimensiones.csv | via=niveles | p
\newcommand{\nSeries}{17}
% Primer mes de la vía en niveles | data/results/05_dimensiones.csv | via=niveles | fecha_inicio
\newcommand{\fechaInicioNiveles}{enero de 2007}
% Primer mes de la vía transformada | data/results/05_dimensiones.csv | via=transformadas | fecha_inicio
\newcommand{\fechaInicioTransformadas}{febrero de 2007}
% Último mes de los datos (ambas vías) | data/results/05_dimensiones.csv | via=niveles | fecha_fin
\newcommand{\fechaFinDatos}{diciembre de 2025}
% Año inicial del período de estudio | data/results/05_dimensiones.csv | via=niveles | fecha_inicio
\newcommand{\anioInicioDatos}{2007}
% Año final del período de estudio | data/results/05_dimensiones.csv | via=niveles | fecha_fin
\newcommand{\anioFinDatos}{2025}
% Inercia (%) de PC1 en niveles | data/results/niveles/05_pca_valores_propios.csv | componente=1 | inercia
\newcommand{\inerciaPCUnoNiveles}{67{,}2}
% Inercia (%) de PC2 en niveles | data/results/niveles/05_pca_valores_propios.csv | componente=2 | inercia
\newcommand{\inerciaPCDosNiveles}{18{,}3}
% Inercia (%) de PC3 en niveles | data/results/niveles/05_pca_valores_propios.csv | componente=3 | inercia
\newcommand{\inerciaPCTresNiveles}{7{,}1}
% Inercia (%) de PC4 en niveles | data/results/niveles/05_pca_valores_propios.csv | componente=4 | inercia
\newcommand{\inerciaPCCuatroNiveles}{2{,}9}
% Inercia (%) de PC1 en transformadas | data/results/05_pca_valores_propios.csv | componente=1 | inercia
\newcommand{\inerciaPCUnoTransformadas}{23{,}3}
% Inercia (%) de PC2 en transformadas | data/results/05_pca_valores_propios.csv | componente=2 | inercia
\newcommand{\inerciaPCDosTransformadas}{16{,}0}
% Inercia (%) de PC3 en transformadas | data/results/05_pca_valores_propios.csv | componente=3 | inercia
\newcommand{\inerciaPCTresTransformadas}{11{,}2}
% Inercia (%) de PC4 en transformadas | data/results/05_pca_valores_propios.csv | componente=4 | inercia
\newcommand{\inerciaPCCuatroTransformadas}{8{,}6}
% Componentes retenidos por Horn en niveles | data/results/niveles/05_pca_retencion.csv | criterio=horn | componentes
\newcommand{\hornNiveles}{2}
% Componentes retenidos por Horn en transformadas | data/results/05_pca_retencion.csv | criterio=horn | componentes
\newcommand{\hornTransformadas}{4}
% Componentes con valor propio > 1 en niveles | data/results/niveles/05_pca_retencion.csv | criterio=kaiser | componentes
\newcommand{\kaiserNiveles}{3}
% Componentes con valor propio > 1 en transformadas | data/results/05_pca_retencion.csv | criterio=kaiser | componentes
\newcommand{\kaiserTransformadas}{6}
% KMO global en niveles | data/results/niveles/05_adecuacion.csv | ambito=global;medida=kmo_global | valor
\newcommand{\kmoGlobalNiveles}{0{,}86}
% KMO global en transformadas | data/results/05_adecuacion.csv | ambito=global;medida=kmo_global | valor
\newcommand{\kmoGlobalTransformadas}{0{,}70}
% Chi cuadrado de Bartlett con n en niveles | data/results/niveles/05_adecuacion.csv | ambito=global;medida=bartlett_chi2_n | valor
\newcommand{\chiBartlettNiveles}{12\,210}
% Chi cuadrado de Bartlett con n en transformadas | data/results/05_adecuacion.csv | ambito=global;medida=bartlett_chi2_n | valor
\newcommand{\chiBartlettTransformadas}{1\,831}
% Grados de libertad de Bartlett (p(p-1)/2) | data/results/niveles/05_adecuacion.csv | ambito=global;medida=bartlett_gl | valor
\newcommand{\glBartlett}{136}
% Chi cuadrado de Bartlett con la mediana de n_ef en niveles | data/results/niveles/05_adecuacion.csv | ambito=global;medida=bartlett_chi2_nef | valor
\newcommand{\chiBartlettNefNiveles}{162}
% p-valor de Bartlett con la mediana de n_ef en niveles | data/results/niveles/05_adecuacion.csv | ambito=global;medida=bartlett_p_nef | valor
\newcommand{\pBartlettNefNiveles}{0{,}06}
% Mediana de n_ef entre pares en niveles | data/results/niveles/05_n_efectivo_resumen.csv | version=original | n_ef_mediana
\newcommand{\nEfMedianaNiveles}{10{,}4}
% Mediana de n_ef entre pares en transformadas | data/results/05_n_efectivo_resumen.csv | version=original | n_ef_mediana
\newcommand{\nEfMedianaTransformadas}{182{,}0}
% Número de condición de R en niveles | data/results/niveles/05_adecuacion.csv | ambito=global;medida=num_condicion | valor
\newcommand{\numCondicionNiveles}{67\,647}
% Número de condición de R en transformadas | data/results/05_adecuacion.csv | ambito=global;medida=num_condicion | valor
\newcommand{\numCondicionTransformadas}{68}
% Correlación de PC1 con el tiempo en niveles | data/results/niveles/05_pca_tendencia_componentes.csv | nombre=PC1 | r_tiempo
\newcommand{\rTiempoPCUnoNiveles}{0{,}993}
% Clase D5 de PC1 en niveles | data/results/niveles/05_pca_tendencia_componentes.csv | nombre=PC1 | clase_d5
\newcommand{\claseDCincoPCUnoNiveles}{DS}
% Clase D5 de PC2 en niveles | data/results/niveles/05_pca_tendencia_componentes.csv | nombre=PC2 | clase_d5
\newcommand{\claseDCincoPCDosNiveles}{DS}
% Clase D5 de PC3 en niveles | data/results/niveles/05_pca_tendencia_componentes.csv | nombre=PC3 | clase_d5
\newcommand{\claseDCincoPCTresNiveles}{DS}
% Clase D5 de PC4 en niveles | data/results/niveles/05_pca_tendencia_componentes.csv | nombre=PC4 | clase_d5
\newcommand{\claseDCincoPCCuatroNiveles}{DS}
% Clase D5 de PC1 en transformadas | data/results/05_pca_tendencia_componentes.csv | nombre=PC1 | clase_d5
\newcommand{\claseDCincoPCUnoTransformadas}{TS}
% Clase D5 de PC2 en transformadas | data/results/05_pca_tendencia_componentes.csv | nombre=PC2 | clase_d5
\newcommand{\claseDCincoPCDosTransformadas}{TS}
% Clase D5 de PC3 en transformadas | data/results/05_pca_tendencia_componentes.csv | nombre=PC3 | clase_d5
\newcommand{\claseDCincoPCTresTransformadas}{TS}
% Clase D5 de PC4 en transformadas | data/results/05_pca_tendencia_componentes.csv | nombre=PC4 | clase_d5
\newcommand{\claseDCincoPCCuatroTransformadas}{indeterminada}
% Series en niveles clasificadas TS (D5) | data/results/niveles/05_clasificacion_d5_resumen.csv | conjunto=series_niveles | TS
\newcommand{\nSeriesTS}{1}
% Series en niveles clasificadas DS (D5) | data/results/niveles/05_clasificacion_d5_resumen.csv | conjunto=series_niveles | DS
\newcommand{\nSeriesDS}{13}
% Series en niveles con clase indeterminada (D5) | data/results/niveles/05_clasificacion_d5_resumen.csv | conjunto=series_niveles | indeterminada
\newcommand{\nSeriesIndeterminadas}{3}
% Proporción del nulo con deriva con PC1 >= observado | data/results/niveles/05_nulo_tendencias.csv | version=con_deriva;medida=inercia_pc1 | proporcion
\newcommand{\nuloConDerivaPCUno}{0{,}902}
% Proporción del nulo con deriva con PC2 >= observado | data/results/niveles/05_nulo_tendencias.csv | version=con_deriva;medida=inercia_pc2 | proporcion
\newcommand{\nuloConDerivaPCDos}{0{,}015}
% Proporción del nulo sin deriva con PC1 >= observado | data/results/niveles/05_nulo_tendencias.csv | version=sin_deriva;medida=inercia_pc1 | proporcion
\newcommand{\nuloSinDerivaPCUno}{0{,}004}
% Proporción del nulo sin deriva con PC2 >= observado | data/results/niveles/05_nulo_tendencias.csv | version=sin_deriva;medida=inercia_pc2 | proporcion
\newcommand{\nuloSinDerivaPCDos}{0{,}592}
% Proporción del nulo sin tendencia con PC1 >= observado | data/results/niveles/05_nulo_tendencias.csv | version=sin_tendencia;medida=inercia_pc1 | proporcion
\newcommand{\nuloSinTendenciaPCUno}{0{,}124}
% Proporción del nulo sin tendencia con PC2 >= observado | data/results/niveles/05_nulo_tendencias.csv | version=sin_tendencia;medida=inercia_pc2 | proporcion
\newcommand{\nuloSinTendenciaPCDos}{0{,}595}
% Mediana de la inercia de PC1 bajo el nulo con deriva | data/results/niveles/05_nulo_tendencias.csv | version=con_deriva;medida=inercia_pc1 | p50
\newcommand{\nuloConDerivaMedianaPCUno}{74{,}6}
% Inercia (%) de PC1 de las series en niveles sin tendencia lineal | data/results/niveles/05_pca_contraste_tendencia.csv | version=niveles_sin_tendencia;medida=inercia;componente=1 | valor
\newcommand{\inerciaPCUnoSinTendencia}{40{,}0}
% KMO de las series en niveles sin tendencia lineal | data/results/niveles/05_pca_contraste_tendencia.csv | version=niveles_sin_tendencia;medida=kmo | valor
\newcommand{\kmoSinTendencia}{0{,}70}
% Percentil 5 de la congruencia bootstrap de PC2 (niveles; L = 24) | data/results/niveles/05_bootstrap_congruencia.csv | L=24;componente=2 | p5
\newcommand{\congruenciaBootPCDosNiveles}{0{,}793}
% Percentil 5 de la congruencia bootstrap de PC4 (niveles; L = 24) | data/results/niveles/05_bootstrap_congruencia.csv | L=24;componente=4 | p5
\newcommand{\congruenciaBootPCCuatroNiveles}{0{,}608}
